\documentclass[arxiv,usenatbib]{iupartex}
% geometry, xcolor, graphicx, amsmath, amssymb, hyperref, natbib and dcolumn
% latex packages are used in the class file, do not redeclare them.

%%%%%%%%%%%%%% Attention to authors %%%%%%%%%%%%%%%%%%%%%%
% template has two output options, referee and arxiv
%%%%%%%%%%%%%%%%%%%%%%%%%%%%%%%%%%%%%%%%%%%%%%%%%%%%%%%%%%
% Don't change the following lines
\usepackage{newtxtext,newtxmath}
\usepackage[T1]{fontenc}

%%%%% Authors - Place Your Packages Here %%%%%
% Only include extra packages if you really need them. Common packages are:
%\usepackage{subcaption}
%\setcitestyle{square,numbers}

%%%% your own commands here %%%%%%%%%%%%%%%
\newcommand{\textbit}[1]{\textbf{\textit{#1}}}
\newcommand{\inspace}[1]{\vspace*{#1}\noindent}

%%%%%%%%%%%%%%%%%%% Title Page %%%%%%%%%%%%%%%%%%%
%\title[Running title]{article title}
\title[Status Report of Year 1 of the CAPIBARA Collaboration]{Status Report of the first year CAPIBARA Collaboration}

% The list of authors, and the short list which is used in the headers.
% If you need two or more lines of authors, add an extra line using \newauthor
%% for corresponding author put \cc after superscripts number
\author[Y1 CAPIBARA Status Report]{%
Joan Alcaide-Núñez$^{1}$ \orcid{0009-0009-2859-0974}
% do not delete the following line
\affsep \\
% List of institutions with line breaks
$^1$Collaboration for the Analysis of Ionic and Photonic Bursts and Radiation, CAPIBARA Collaboration, Barcelona, Spain
}

%corresponding author name and email
\corres{J. Alcaide-Núñez}{capibara-mission@outlook.com}

% Enter the current year, for the copyright statements etc.
\pubyear{2025}

%%%%%%%% Authors are not allowed to edit %%%%%%%%%
% Don't change the lines up to \maketitle
\doiheader{XXXXXXX/PAR.20XX.00000}
\date{https://www.overleaf.com/project/68120e07791acb82c6236123#
	\pSubmit{00.00.0000} 
	\pRevReq{00.00.0000}
	\pLastRevRec{00.00.0000}
	\pAccept{00.00.0000}
	\pPubOnl{00.00.0000}
}
\volume{0}
\volnumber{1}

\begin{document}
\label{firstpage}
\pagerange{\pageref*{firstpage}--\pageref*{lastpage}}
\maketitle

% Abstract of the paper
% \begin{abstract}
% This is a simple template for authors to write new PAR papers. The abstract should briefly describe the aims, methods, and main results of the paper. It should be a single paragraph, not more than 300 words. No references should appear in the abstract.
% \end{abstract}

% Select between one and six entries from the list of approved keywords.
% Don't make up new ones.
% \begin{keywords}
% keyword1 -- keyword2 -- keyword3
% \end{keywords}

%%%%%%%%%%%%%%%%%%%%%%%%%%%%%%%%%%%%%%%%%%%%%%%%%%

%%%%%%%%%%%%%%%%% BODY OF PAPER %%%%%%%%%%%%%%%%%%

\vspace{0.5cm}
\section{Introduction}
Approaching the first year since the creation of the Collaboration for the Analysis of Ionic and Photonic Bursts and Radiation (CAPIBARA Collaboration), this letter aims to provide an overview of the project's current status and serve as a guidelines for coming work. This document is primarily written for members of the Collaboration and not specifically for partnership

This document is primarily written for the members of the Collaboration and not for other partnership institutions, companies, or entities.

\section{Progress in the First Year}
Since the inception of the CAPIBARA Collaboration in summer 2024 (July), we have been defining the group and our goals. Towards the end of that year, we submitted our first application to the \href{https://spark-program.pldspace.com}{SPARK Program} by \href{https://pldspace.com}{PLD Space} elaborating on our vision in a \href{https://capibara3.github.io/files/preliminar_proposal.pdf}{preliminary report}, to exploring the high-energy cosmos in both its ionic and photonic aspects.

In early 2025, we participated in the matchmaking session proposed by PLD Space, where we presented our project and aims. From that matchmaking session, we were able to learn about other initiatives and companies, and later contact them for potential partnerships and support. Addressing our goals more realistically, we submitted a partnership report to the SPARK program outlining each of our goals and technical details, as well as a description of our partnerships:
\begin{itemize}
    \item LabOSat (Detector technology)
    \item TinyGS (Telemetry)
    \item TeideSat (Outreach)
    \item IFAE (Detector technology, only if we were accepted by SPARK)
\end{itemize}

However, the change in orbit by the SPARK Program, the duality of our goals and our technical possibilities as high school students imposed the need to divide our objectives into separate missions. As the cosmic ray detector technology is more mature in our group and its technical requirements are less, we focused on the ionic goal, i.e., the CAPIBARA Cosmic Rays Mission (CAPICR\footnote{This abbreviation is the one used in GitHub, but we still need to find a name for it.}). CAPICR consists of a Time-of-Flight (ToF) detector and potentially a charge detector, providing information on the velocity and charge of the particles trapped on Earth's orbit. The prioritization of this mission comes from the fact that some members of the group were already familiar with the detectors and the needs (in comparison with a high-energy observatory) are much simpler. Although they don't have the same goals, we see CAPICR as a pathfinder for CAPIGX, since it will enable us to gain knowledge and expertise into the aerospace field, leveraging our position to bring CAPIGX forward.

At this point, we also contacted Orbital Boost Aerospace, a satellite booster company which expects to launch the OBA FARADAY satellite, as a proof of concept mission, with the SPARK Program. Their 6U CubeSat has 1U free space, which would fit our cosmic ray detector, thus we established a partnership with OBA, who submitted an updated report to PLD Space, this time with CAPICR as part of the OBA FARADAY mission. Note that there is already another hosted instrument from \href{https://www.lnaspace.com}{LNA Space}

\section{Current Tasks on CAPICR}
Our recent efforts have focused on completing a detector design a cosmic ray detector, tailoring each piece, as well as searching for options to do a crowdfunding campaign for the funding of this instrument. Nevertheless, progress was slowed by the selectividad exams, and we hope to advance more in the future.

\textbf{Detector Design and Construction}\newline
Before the start of the academic year, we should have a proper design for the cosmic ray detector, in cohesion with the components and costs estimated in the \href{https://github.com/CAPIBARA3/capicr-detector-design/blob/main/Cosmic_Ray_Detector_Materials.md}{components list} (which should be already finished). We will have to comply and keep up with the progress of the OBA FARADAY mission, and be ready for integration when they reach that step.

\textbf{Financial Status: Crowdfunding Campaign}\newline
Regarding our financial status, we currently lack support from any entity of company in this aspect. Therefore, the CAPIBARA Collaboration must fund itself. We are currently building a crowdfunding campaign on \href{https://whydonate.com}{WhyDonate}. For this campaign, we will require audiovisual materials, social media \footnote{@capibara\_collaboration on \href{https://www.instagram.com/capibara_collaboration/}{Instagram}, CAPIBARA Collaboration on \href{https://www.linkedin.com/company/107744838/admin/dashboard/}{LinkedIn}} activity (i.e. starting accounts on behalf of the collaboration), and reaching out to everyone you know of. We are currently starting the social media and creating visual content for the campaign, expected to launch later this year. In the face of lack of funding we would discuss extending the crowdfunding and/or contacting foundations and institutions for support.

\section{Plans for CAPIGX}
As stated above, CAPIBARA has dual goals, while one of them is to put a cosmic ray detector in space (the ionic goal), we have not forgotten our aim to study the electromagnetic high-energy Universe. CAPIBARA's Gamma-/X-ray Mission (CAPIGX) is the second mission in which we separated our initial vision. Focusing on the electromagnetic side of the high-energy Universe, our goal is to provide accurate gamma-ray and X-ray transient observations supporting bigger telescopes. The role of CAPIGX is to provide a follow-up and discovering observatory for high-energy transients during the multi-messenger era (2030s), at the same time as the next generation gravitational wave measurements (Cosmic Explorer, Einstein Telescope, LISA) and high-profile X-ray and Gamma-ray missions (NewAthena, THESEUS).

CAPIGX will be sensitive in the $1$ keV - $1$ MeV energy range (with a possible extension to 5 MeV if possible with the same detector technology). The extension to X-ray enabling high-$z$ transient detection, as these have been redshifted and therefore tend to be visible in the X-ray band. Additionally, the CAPIGX will consist of 3 different CubeSat detectors, allowing for intensity interferometry, a technique in which we use the different signal times of arrival at the CubeSats to provide better source localization accuracy, exponentially enhancing the constrains. This is key, for instance, for the discovery of host galaxies of GW events. CAPIGX is currently only on paper, but we are planning big steps for the future. Our next efforts and a summary of the timeline of the mission is shown below. Key to this mission is to develop a robust, realistic, and appropriate application to FlyYourSatellite! taking into account the status of the fields, the planned mission, the program's requirements and expectations, our ambitions and the timeline.
\begin{itemize}
    \item 2025: Feasibility Study and Writing of Paper Justifying and Outlining the Characteristics of the Mission
    \item 2026: Detailed Engineering Design and Partnership Building with Manufacturers
    \item 2027: Apply to \href{https://www.esa.int/Education/CubeSats_-_Fly_Your_Satellite}{FlyYourSatellite!} Program by ESA\footnote{Note that this is subjects to the call for proposals of ESA and hence possible changes are probable.}
    \item 2028-2029+: Finish Development, Construction, Integration
    \item 2030+: Launch Window (expected for 2032-2035)
    \item $\sim$2040: End of Mission\footnote{Since CAPIGX is a CubeSat constellation it can be easily upgraded and expanded depending on the success and needs of the field.}
\end{itemize}


\section{Outlook}
In a few words, the CAPIBARA Collaboration has been quite productive in its first year. We are working on acquiring the materials for the construction of the CAPICR mission right now and expecting to have begun the built before the end of 2025. Furthermore, we are in contact with OBA Space and PLD Space awaiting for a green light for our joint launch opportunity. Regarding CAPIGX, we have a general outline of the progress we need to do and are expecting to publish a preliminary study (\textit{Alcaide-Núñez in prep.}), in which we will outline the status of the field and the observational gap, as well as the details related to instruments and challenges such as intensity interferometry.


\section{Acknowledgments}
The work described and presented here has been developed by a \textit{collaboration} of students. Additionally, the CAPIBARA Collaboration appreciates the opportunities given by our partnerships, the Spark Program, and all the support received.
\vspace{0.5cm}
The author(s) of this report would appreciate comments and suggestions, please reach out to \href{mailto:capibaracollaboration@outlook.com}{capibaracollaboration@outlook.com}.

\section*{Notes}
We use GitHub as our main information hub, hosting documents, files, code and planning information (tasks, timelines). The most important sites are our \href{https://github.com/orgs/CAPIBARA3/repositories}{repositories overview} and the project's \href{https://github.com/orgs/CAPIBARA3/projects/1/views/1}{timeline view}, which also has a tab with unassigned tasks. An overview of current repositories is here:
\begin{itemize}
    \item \textbf{\href{https://github.com/capibara3/.github}{.github}:} Public Profile Repository
    \item \textbf{\href{https://github.com/capibara3/.github-private}{.github-private}:} Internal Repository for templates, OnBoarding Members, Tools Tutorials, ...
    \item \textbf{\href{https://github.com/capibara3/capicr-software}{capicr-software}:} Repository for Onboard and ground software of CAPICR
    \item \textbf{\href{https://github.com/capibara3/capicr-detector-design}{capicr-detector-design}:} Repository for the Engineering and Design of the CAPICR Instrument
    \item \textbf{\href{https://github.com/capibara3/capigx-obs-stats}{capigx-obs-stats}:} Repository for Statistics About the Status of X-/Gamma-Ray Observatories
    \item \textbf{\href{https://github.com/capibara3/capigx-data-analysis}{capigx-data-analysis}:} Repository for CAPIGX Data Production and Analysis Simulations
    \item \textbf{\href{https://github.com/capibara3/cosmology}{cosmology}:} Repository for Simulations on Cosmology with GRBs at high-z and Different Models
    \item \textbf{\href{https://github.com/capibara3/capibara3.github.io}{capibara3.github.io}:} Repository for the Website
\end{itemize}




\end{document}
